\documentclass[conference]{IEEEtran}
\usepackage[T1]{fontenc}
\usepackage[utf8]{inputenc}
\usepackage{amsmath,amssymb,booktabs,graphicx,url,array}
\usepackage[hidelinks]{hyperref}
\title{Verifying AI-Generated Legal Citations Against Indian Case Law: An Evidence-First Integrity Pipeline}
\author{\IEEEauthorblockN{Afjal Hussein, Ayush Kumar, Arush Sharma}
\IEEEauthorblockA{S24BCAU0171, S24BCAU0103, S24BCAU0161\\
Department of Computer Science Engineering \& Technology\\Bennett University, India\\
CSET485---AI and Society, Phase 2\\Team ID: not supplied in the source material}}
\begin{document}
\maketitle
% Generated by scripts/build_report.py from the recorded experiment.
\newcommand{\CorpusCount}{385}
\newcommand{\ParagraphCount}{13,456}
\newcommand{\RequestedCount}{2,000 (1,873 selected keys)}
\newcommand{\DevJudgments}{115}
\newcommand{\TestJudgments}{270}
\newcommand{\TestCount}{2,421}
\newcommand{\CpuName}{Intel Core i5-12450HX}
\newcommand{\RamSize}{15.78 GiB}
\newcommand{\PythonVersion}{3.11.9}
\newcommand{\TorchVersion}{2.6.0+cu124}
\newcommand{\FullFone}{0.986}
\newcommand{\ExistFone}{0.251}
\newcommand{\ComparisonRows}{Citation existence & 0.335 & 1.000 & 0.144 & 0.251\\
Literal quote & 0.356 & 0.687 & 0.315 & 0.432\\
Normalized quote & 0.446 & 1.000 & 0.287 & 0.446\\
Full verifier & 0.979 & 1.000 & 0.973 & 0.986\\}
\newcommand{\RetrievalRows}{BM25 test R@1 / R@5 & 1.000 / 1.000\\
BM25 test MRR & 1.000\\
MiniLM test R@1 / R@5 & 0.948 / 0.999\\
MiniLM test MRR & 0.972\\
Test retrieval queries & 804\\
Model revision prefix & 1110a243fdf4\\
Token limit / dimensions & 256 / 384\\
Selected inference batch & 16\\
Paragraph indexing time & 19.97 s\\}
\newcommand{\TuningRows}{0.8 & 0.25 & 0.9971 & 0.9985\\
0.8 & 0.75 & 1.0000 & 1.0000\\
1.5 & 0.25 & 0.9942 & 0.9971\\
1.5 & 0.75 & 0.9971 & 0.9985\\
2.0 & 0.25 & 0.9942 & 0.9971\\
2.0 & 0.75 & 0.9942 & 0.9971\\}

\begin{abstract}
Generative legal research can produce plausible citations while attaching inaccurate quotations, incorrect court information, or unsupported pinpoints. This project implements a local, evidence-first verifier for a bounded collection of Indian Supreme Court judgments. The system resolves metadata-backed citation aliases, compares court and bench fields, checks original paragraph identifiers, searches for contiguous normalized quotations, and exposes nearby qualifying language for human review. It separates textual consistency from legal entailment and treats an absent local citation as unresolved rather than fabricated. A reproducible acquisition process obtains public judgment PDFs and metadata, records source URLs and SHA-256 hashes, and rejects unusable extractions. Experiments use \CorpusCount{} retained judgments and judgment-disjoint development and test partitions. Controlled mutations yield \TestCount{} held-out verification requests. On this semi-synthetic triage task, the complete pipeline obtains an F1 score of \FullFone{}, compared with \ExistFone{} for citation-existence checking. These measurements concern automatically constructed discrepancies, not the prevalence of hallucinations in actual litigation. A separate retrieval experiment compares lexical ranking with pretrained sentence embeddings on the available GPU. The implementation, test suite, local browser demonstration, per-request predictions, and editable report source are delivered together. The contribution is an inspectable workflow that reduces repetitive checking without representing a successful lookup as a legal conclusion.
\end{abstract}
\begin{IEEEkeywords}
legal citation verification, Indian case law, information retrieval, quotation integrity, hallucination detection, responsible AI
\end{IEEEkeywords}

\section{Introduction}
Legal citations make actionable claims about sources. A citation identifies an authority, a pinpoint directs the reader to particular text, and a quotation represents what that authority said. These claims can fail independently. A judgment may exist but have been issued by a different court. Its text may contain a similar sentence while differing in a negation, amount, or condition. A passage may be reproduced accurately while the proposition attached to it goes beyond what the judgment establishes. A useful verification tool must therefore avoid collapsing all of these questions into one similarity score.

Large language models make this separation important. Research on legal hallucinations documents failures on verifiable questions about real cases, while audits of commercial legal research systems test reliability in products intended for professional use. These studies motivate independent checking \cite{ref1,ref2}, but their United States results cannot be transferred to Indian practice or to this project's dataset. Retrieval-augmented generation conditions a model on sources; it does not guarantee that every generated citation or quotation is correct. A verifier should show source evidence and uncertainty rather than reproduce that guarantee under a different name \cite{ref3,ref4}.

For Indian case law, citation aliases and document structure create practical difficulties. A judgment may have neutral and SCR citations, while a user's reference uses a commercial reporter absent from the available metadata. PDFs include headnotes, page numbers, cited passages, and sometimes more than one opinion. Treating every numbered fragment as a judgment paragraph would create misleading evidence. The engineering problem addressed here is narrower than a general legal assistant: identify the claimed authority, inspect available metadata and extracted text, and state which mechanical claims can be checked.

\subsection{Problem and Measurable Objectives}
The input is a citation string with a quotation, a proposition, or both; court, bench size, and paragraph number are optional. The output is an evidence-bearing report with one of four states: consistent, inconsistent, review, or unresolved. Consistent means the submitted mechanical claims match the local extraction. It does not certify the proposition, the passage's status as a holding, or the authority's continuing validity.

The objectives are to (1) ingest a bounded public corpus with document provenance; (2) resolve only metadata-supported aliases; (3) detect controlled quotation and metadata changes while preserving numbers and negations; (4) return original paragraph identifiers and source-linked evidence; and (5) compare the pipeline with three simpler baselines on identical held-out cases. Extraction retention, uncertainty outcomes, retrieval quality, and measured latency complement classification metrics because accuracy alone can conceal weak source coverage.

\subsection{Societal Relevance and Contributions}
An offline, inspectable tool can help legal-aid workers, students, and advocates perform repetitive source checks without sending draft text to an external model. Public court records provide an alternative to the inaccurate claim that verification always requires a paid database. Nevertheless, a limited English-language Supreme Court sample cannot cover High Courts, regional languages, every reporter, or historical treatment. Unequal corpus coverage can create unequal usefulness: a missing result must not be presented as a finding of fabrication.

The contribution is an implemented verification workflow, not doctrinal reasoning. It includes provenance-aware acquisition, conservative paragraph extraction, deterministic checks, a local CLI and browser interface, and controlled experiments. Phase 1 discussed wrong-voice errors and subsequent treatment as part of the broader problem; ratio/obiter classification, overruling detection, and autonomous drafting remain outside the committed implementation. No commercial system is claimed to lack a capability merely because its implementation is unpublished. The project prioritizes a reproducible demonstration of the agreed subset over unverifiable claims of comprehensive legal safety.

\section{Literature Survey}
\subsection{Hallucination Audits and Source Grounding}
Dahl \emph{et al.} pose specific, verifiable questions about United States cases and compare model responses with known legal facts \cite{ref1}. Their taxonomy distinguishes several kinds of legal hallucination and their experiments show why fluent answers cannot be treated as reliable evidence. The contribution to this project is a motivation for independently checking identifiable claims. Its limitation is transfer: a United States audit neither establishes the error distribution in Indian filings nor measures this detector's performance. We consequently do not use its hallucination rates as target accuracy figures.

Magesh \emph{et al.} audit leading commercial AI legal research tools rather than general-purpose models \cite{ref2}. The study tests vendor reliability claims using legal research queries and evaluates generated responses. It motivates checking even systems that retrieve legal sources. However, these are end-to-end product results, while our controlled experiment concerns predefined mechanical changes to source-linked requests. Directly comparing those rates with our F1 score would compare different populations, labels, and tasks. We neither reproduce proprietary products nor claim superiority to their legal reasoning.

Lewis \emph{et al.} combine a pretrained generator with retrieval from nonparametric memory for knowledge-intensive tasks \cite{ref3}. This supports the architectural value of access to external evidence, but retrieval and verification remain distinct. The presence of a retrieved passage does not establish that the generated quotation reproduces it accurately. Gao \emph{et al.} introduce ALCE to evaluate answers with citations, including their correctness and completeness \cite{ref4}. Their work demonstrates why attaching a citation is weaker than supporting the full claim. General question-answering corpora do not provide the Indian court, reporter-alias, and numbered-paragraph annotations needed here. Our implementation therefore avoids claiming that a text lookup solves citation entailment.

\subsection{Fact Verification and Evaluation Limits}
FEVER introduces a large claim-verification dataset using Wikipedia evidence and supported, refuted, or insufficient-information labels \cite{ref5}. Its distinction between a negative verdict and insufficient evidence informs this project's unresolved state. Nevertheless, evidence selection in encyclopedic text is not a substitute for identifying the deciding court's voice or the scope of a legal holding. SciFact pairs expert-written scientific claims with abstracts and supporting or refuting evidence \cite{ref6}. The expert-annotation component is particularly relevant: automatically generating legal error labels is convenient, but cannot replace independent legal judgment. Neither benchmark is used as a surrogate Indian evaluation set.

FactCC studies factual consistency between source documents and abstractive summaries using synthetic transformations \cite{ref7}. This supplies a useful precedent for controlled perturbations while exposing their limitation: good performance on artificial transformations does not ensure detection of naturally occurring errors. SelfCheckGPT instead uses consistency across sampled model generations to detect hallucinations without an external reference collection \cite{ref8}. Such sampling provides a different reliability signal from this project's explicit judgment evidence. When an original source is available, directly inspecting the claimed words offers a simpler, explainable check; it still leaves the proposition's legal meaning unresolved.

TRUE standardizes evaluation of factual-consistency metrics across annotated datasets \cite{ref19}. Its emphasis on example-level assessment and variation between tasks supports preserving per-case predictions, confusion counts, and category breakdowns. A single aggregate number can conceal different failure costs. In legal checking, a false accusation that a case is invented and an unnecessary human-review warning are not interchangeable outcomes. We therefore report triage and four-state agreement separately rather than relabelling every warning as a proven false citation.

\subsection{Legal Language Resources}
LegalBench organizes diverse legal reasoning tasks through a collaborative benchmark \cite{ref9}. It supports task-specific capability evaluation, but a broad legal benchmark does not supply an operational Indian citation-indexing pipeline. LexGLUE standardizes multiple English legal language-understanding datasets and compares generic with legal-oriented models \cite{ref10}. Its findings support domain-sensitive evaluation. Document classification performance, however, says little about exact quotation fidelity, a missing paragraph identifier, or an authority's subsequent treatment.

Malik \emph{et al.} introduce ILDC for court judgment prediction and explanation using Indian Supreme Court documents \cite{ref11}. The source domain is closely relevant, and the work demonstrates that Indian judgments can support reproducible NLP research. Its target remains different: predicting an outcome from a judgment is not verifying a quotation supplied by a reader. This project therefore uses an attributed public source and defines a separate integrity task instead of presenting an outcome-prediction score as evidence of citation accuracy.

BERT learns contextual language representations through bidirectional pretraining \cite{ref12}. LEGAL-BERT investigates pretraining and adaptation for legal text \cite{ref14}, showing why domain mismatch deserves attention. Neither pretraining objective certifies legal support. A legal-domain encoder could improve candidate retrieval, but its score would still require task-specific evaluation and would not replace source-level metadata checks. We use deterministic checks for the mechanical requirements and reserve pretrained representations for a separately reported retrieval experiment.

\subsection{Retrieval Algorithms and the Remaining Gap}
Sentence-BERT adapts transformer representations for efficient sentence-level similarity search \cite{ref13}. The Transformer architecture supplies the underlying contextual attention mechanism \cite{ref18}. These approaches make vector comparison practical, but generic sentence similarity can reward paraphrases that differ in a legally decisive detail. Dense Passage Retrieval learns question--passage representations for open-domain question answering \cite{ref16}; its relevance criterion and candidate collection differ from within-judgment quotation localization. The current experiment therefore compares candidate ranks without treating dense similarity as an integrity verdict.

Robertson and Zaragoza explain probabilistic relevance ranking, including BM25's term-frequency saturation and document-length normalization \cite{ref15}. BM25 is attractive here because a quoted fragment usually shares vocabulary with its source and the score is inexpensive to compute. BEIR evaluates retrieval methods across heterogeneous zero-shot tasks and shows that effectiveness depends on domain and dataset \cite{ref17}. Together these results motivate a measured lexical-versus-dense comparison rather than assuming that GPU use necessarily improves retrieval.

The gap addressed is an inspectable workflow connecting Indian citation aliases, original paragraph numbers, quotation fidelity, mechanical metadata checks, and uncertainty reporting. We do not claim this is the first or only such system. The literature instead supports a bounded engineering conclusion: evidence presentation, explicit task definitions, and abstention remain necessary even with strong language models. The project contributes a runnable implementation and transparent measurement of that subset. Independently annotated natural errors, voice attribution, legal entailment, and subsequent-treatment analysis remain separate research problems.


\section{Proposed Methodology}
\subsection{Architecture and Data Flow}
The pipeline comprises acquisition, extraction, indexing, checking, and evidence presentation (Fig.~\ref{fig:pipeline}). Acquisition and experiments are separate commands; an ordinary request never silently downloads a judgment or calls an LLM. The same Python function drives the command-line and browser interfaces, so demonstration and measurement exercise the same engine.

\begin{figure}[t]\centering
\setlength{\fboxsep}{6pt}
\fbox{\parbox{0.83\columnwidth}{\centering Public registry metadata + judgment PDFs\\$\downarrow$\\Bounded downloads, hashes, original paragraph extraction\\$\downarrow$\\Metadata-backed citation index\\$\downarrow$\\Court / bench / pinpoint / quote checks\\$\downarrow$\\Evidence paragraphs + uncertainty-aware findings}}
\caption{Implemented workflow. Optional dense retrieval ranks candidates; it never overrides a failed quotation check.}\label{fig:pipeline}
\end{figure}

Each corpus row contains a stable ID, case name, decision date, court, known bench size, citation aliases, numbered paragraphs, source and metadata URLs, and PDF hash. Paragraphs retain their original number, text, and source page. The loader rejects malformed records, duplicate IDs, duplicate paragraph identifiers, invalid dates, non-HTTP source links, and invalid metadata types. Aliases normalizing to the same identifier within one record are deduplicated. When an alias identifies different records, the engine abstains rather than choosing arbitrarily.

\subsection{Acquisition and Paragraph Extraction}
The source is the Indian Supreme Court Judgments collection listed by the AWS Open Data Registry under CC-BY-4.0 \cite{dataset}. The collector selects a deterministic bounded set of metadata keys, downloads English PDFs with retries and timeouts, and caches raw files. The manifest preserves requested and retained counts, errors, retrieval time, source location, selection procedure, and attribution. Downloaded documents are data, not an execution environment. Source-data licensing does not grant reuse rights for project code.

Extraction seeks explicit paragraph markers in the judgment body rather than converting page positions into paragraph numbers. Headnotes and pagination require filtering. Repeated numbering can indicate quoted material, lists, or another opinion, so ambiguous material is rejected or conservatively truncated rather than renumbered. Successful extraction does not prove that every paragraph was captured. A missing pinpoint in an incomplete extraction therefore produces an unresolved finding, not a confident assertion that the paragraph never existed. These precautions trade recall for a more honest evidence trail.

The data card documents retained coverage and known biases. Acquisition targets and usable judgment counts differ: every rejected PDF remains accounted for in the manifest. The corpus is a research slice, not an authoritative citator. A professional user should inspect the original judicial record linked by the tool. The automatic extraction has not received a complete expert paragraph-by-paragraph audit, and controlled experiments cannot substitute for that audit.

\subsection{Citation Resolution and Metadata Checks}
Canonicalization standardizes case, whitespace, neutral-citation spacing, and punctuation in reporter abbreviations. SCC, AIR, SCR, SCALE, and INSC strings can be represented, but an alias resolves only when actually present in corpus metadata. Parsing SCC does not imply that this SCR/neutral-focused source supplies SCC mappings. The engine does not infer aliases from case-name similarity: common party names and unrelated reporter numbers can yield false matches.

Let $a$ be the canonical citation and $I(a)$ its indexed record set. If $|I(a)|=0$, the result is unresolved; if $|I(a)|>1$, identity is ambiguous. Checks proceed only when $|I(a)|=1$. Known court aliases such as SC are normalized. A mismatched requested court or known bench size generates a discrepancy; missing bench metadata produces review. Bench size alone is not used to infer binding force, which also depends on jurisdiction, later treatment, and the legal issue.

\subsection{Quotation Matching and Candidate Retrieval}
Text normalization applies Unicode compatibility normalization, case folding, English tokenization, and whitespace normalization while preserving numeric and negation tokens. A quotation must occur as an ordered, contiguous sequence with token boundaries. Consequently, ``restricted'' cannot match inside ``unrestricted'', while deleting ``not'' or changing an amount prevents a match. Punctuation is normalized, so the result concerns normalized text rather than typographic identity with the PDF.

Let $N(q)$ and $N(p)$ be normalized quotation and paragraph token sequences. The check is
\begin{equation}
V(q,p)=\mathbf{1}[N(q)\text{ is a contiguous subsequence of }N(p)].
\end{equation}
If a pinpoint $n$ is supplied, the engine checks that the quotation occurs there, distinguishing a match elsewhere from a missing or unverified number. A quote can occur more than once; a valid occurrence at the requested pinpoint is sufficient. Candidate retrieval cannot create a false verified pinpoint.

When no contiguous match exists, token coverage can expose a nearby candidate:
\begin{equation}
C(q,p)=\frac{|T(q)\cap T(p)|}{\max(1,|T(q)|)},
\end{equation}
where $T(x)$ is the set of normalized tokens. The fixed candidate-display threshold is 0.3; it affects evidence presentation, not the exact-match verdict, and is not claimed to be tuned. Evidence is labelled verified, candidate, context, or retrieved. A high-overlap candidate cannot silently become a successful quotation check. A separate proposition always requires human review, even when its accompanying quotation matches. There is no learned entailment classifier or task-specific model training.

\subsection{Selective Quotation and Decision Semantics}
A heuristic inspects text immediately following a uniquely matched quote and the beginning of the next numbered paragraph for qualification cues such as ``however'', ``unless'', or ``subject to''. Available original context is shown. The finding is called context review because a cue may be irrelevant, and important qualifications may contain no cue. It does not establish deliberate misrepresentation. Multi-paragraph spans, nested quotations, attribution to different speakers, and doctrinal scope remain difficult cases.

Inconsistent means a detectable mechanical discrepancy; review means a proposition, context cue, missing metadata, or extraction uncertainty requires attention. Unresolved means citation identity could not be established uniquely. Consistent means submitted checks passed within the extraction. Stable finding codes and source evidence accompany the state; limitations are part of the response rather than hidden in documentation.

The local server binds to loopback by default, validates input, and renders user text as text rather than executable markup. Ordinary verification makes no model API calls. Public judgments may contain sensitive facts, so public availability is not unlimited ethical permission to aggregate personal information. The project distributes attributed research data, excludes raw download caches and original course PDFs from Git, and recommends source inspection rather than reliance on a colored status alone.

\section{Experimental Results}
\subsection{Dataset and Computational Environment}
Table~\ref{tab:data} records the actual run. The manifest and corpus hash identify the data. The requested limit of 2,000 exceeded the 1,873 available selected keys in 2022--2023; 385 records were retained (20.6\%), and 1,488 failures remain logged. Most rejected records had no unambiguous body-paragraph extraction; two PDFs returned HTTP 404. The retained slice contains 62 judgments from 2022 and 323 from 2023. This imbalance reflects extraction and source availability, not a representative sample of Indian disputes. Judgment IDs are split before request variants are constructed, preventing variants of a judgment from appearing in both development and test sets. There is no gradient-based training. Optimizer, learning rate, epochs, and training batch size are not applicable; encoder batch size is an inference setting.

\begin{table}[t]\centering\caption{Corpus and experimental environment.}\label{tab:data}
\small\begin{tabular}{>{\raggedright\arraybackslash}p{0.42\columnwidth}>{\raggedright\arraybackslash}p{0.48\columnwidth}}\toprule
Item & Recorded value\\\midrule
Acquisition requests & \RequestedCount{}\\
Retained judgments & \CorpusCount{}\\
Extracted paragraphs & \ParagraphCount{}\\
Development / test judgments & \DevJudgments{} / \TestJudgments{}\\
Held-out verification requests & \TestCount{}\\
CPU & \CpuName{}\\
RAM & \RamSize{}\\
GPU & NVIDIA RTX 3050 Laptop, 6 GB\\
Python / PyTorch & \PythonVersion{} / \TorchVersion{}\\\bottomrule
\end{tabular}\end{table}

The benchmark is semi-synthetic. Correct requests come from real extracted text; controlled variants alter metadata, a quote, a pinpoint, or local citation availability. Unknown bench sizes cannot provide reliable automatic targets and are excluded from bench mutations. Missing pinpoints in incomplete extraction remain uncertain rather than being relabelled as proven nonexistent paragraphs. Context-omission labels describe a construction rule, not expert legal judgments.

This makes the experiment reproducible but creates a validity limitation: the generator and detector share a mechanical definition of error. High scores show that implemented checks respond to controlled conditions, not that they generalize to naturally occurring hallucinations, accurate quotations paired with false propositions, or adversarial filings. An external benchmark must be independently annotated by qualified reviewers with disagreement resolution. The corpus's acquisition failures and extraction uncertainty are therefore reported alongside classification scores.

\subsection{Metrics and Baselines}
For triage, the positive label denotes a request constructed to need attention, not a declaration of fabrication. Accuracy is $(TP+TN)/n$, precision $P=TP/(TP+FP)$, recall $R=TP/(TP+FN)$, and $F_1=2PR/(P+R)$. Here $TP$, $TN$, $FP$, and $FN$ are standard confusion counts relative to this label; zero denominators yield zero. Warnings and discrepancies can both trigger triage, while their separate output states remain available in per-case predictions.

Three baselines use identical held-out requests. Citation-only flags absent identifiers but ignores text and metadata. Literal-quote checking adds direct string matching and is sensitive to harmless formatting changes. Normalized-quote checking tolerates formatting but omits the full metadata and context checks. The complete pipeline adds those checks and uncertainty handling. These are transparent engineering ablations, not reimplementations of commercial legal products or state-of-the-art doctrinal reasoning systems.

\begin{table}[t]\centering\caption{Held-out triage on identical controlled requests.}\label{tab:comparison}
\small\begin{tabular}{lrrrr}\toprule
Method & Acc. & Prec. & Rec. & $F_1$\\\midrule
\ComparisonRows
\bottomrule\end{tabular}\end{table}

The full pipeline obtains \FullFone{} F1 on \TestCount{} requests. Its difference from citation-only checking reflects additional information rather than a newly learned legal representation. A real identifier remains unchanged in wrong-court, wrong-bench, and corrupted-quote requests, leaving an existence-only method unable to flag them. Literal matching also treats formatting variation as an error. Normalization addresses that failure mode, and metadata/pinpoint checks cover claims that quotation matching alone cannot settle.

\subsection{Development Tuning and GPU Retrieval}
The paragraph-localization experiment compares BM25 with pretrained sentence embeddings. Candidates belong to the resolved judgment; a query is an automatically derived source-paragraph fragment. This is within-document localization, not open-domain legal question answering. Candidate sets and relevance labels are identical for all methods. Recall at rank one and five measures return of the original paragraph; repeated wording can create additional relevant paragraphs not recognized by the automatic label.

BM25 computes
\begin{equation}
\sum_{t\in q}\operatorname{IDF}(t)\frac{f(t,p)(k_1+1)}{f(t,p)+k_1(1-b+b|p|/\overline L)},
\end{equation}
where $f(t,p)$ is term frequency, $|p|$ paragraph length, $\overline L$ mean candidate length, $k_1$ saturation, and $b$ length normalization. Dense retrieval compares normalized embeddings with cosine similarity. Long paragraphs face a recorded encoder token limit, unlike lexical ranking. Inference batching affects memory and throughput rather than relevance labels.

\begin{table}[t]\centering\caption{Recorded tuning and retrieval measurements.}\label{tab:retrieval}
\small\begin{tabular}{p{0.55\columnwidth}p{0.35\columnwidth}}\toprule
Measurement & Value\\\midrule
\RetrievalRows
\bottomrule\end{tabular}\end{table}

The BM25 grid in Table~\ref{tab:tuning} uses 344 development queries from 115 judgments. It selects $k_1=0.8$ and $b=0.75$ by maximum development MRR, then recall at one and deterministic parameter tie-breaking. The selection is frozen before evaluating 804 test queries. Perfect BM25 retrieval on this task is unsurprising because queries reuse source tokens and candidates are limited to one judgment; it does not establish semantic search quality.

\begin{table}[t]\centering\caption{BM25 development-only grid; 344 source-fragment queries.}\label{tab:tuning}
\small\begin{tabular}{rrrr}\toprule
$k_1$ & $b$ & Recall@1 & MRR\\\midrule
\TuningRows
\bottomrule\end{tabular}\end{table}

The GPU comparator uses all-MiniLM-L6-v2 with 384-dimensional normalized embeddings and a 256-token limit. Batch sizes 8, 16, and 32 took 0.334, 0.273, and 0.274 seconds on 128 development paragraphs after warm-up; size 16 was selected for throughput. These single-run timing differences are not statistically conclusive, and no accuracy tuning or fine-tuning of weights occurred. Each corpus paragraph was embedded once; scores were computed within the resolved judgment rather than using a global query--document matrix. The encoder was pinned to the revision recorded in the results and loaded from safe tensor weights without remote model code.

MiniLM returned the source paragraph first for 94.8\% of test fragments, below BM25's 100\%, while recall at five reached 99.9\%. Generic semantic representations and truncation can hurt precise lexical localization; the experiment does not isolate their individual effects. The GPU indexed 13,456 paragraphs in 19.97 seconds. Retrieval inference is optional: these embeddings do not alter the core verifier or its quoted-text decisions.

\subsection{Error Analysis and Reproducibility}
The full verifier's 51 missed attention cases all came from the 265 selective-quote constructions: 214 received the expected review state (80.8\%). All 540 clean/whitespace requests were accepted, giving $TP=1{,}830$, $TN=540$, $FP=0$, and $FN=51$ for triage. Controlled missing-pinpoint cases received review, and absent local identifiers received unresolved. Four-state agreement was also 97.89\%, but these state distinctions should still be retained: a warning is not a proven legal falsehood. The benchmark includes 270 cases in each common category, 269 wrong-bench cases, 267 wrong-pinpoint cases, and 265 selective-quote cases. The attention-positive class is 77.7\% of the test set, so the class distribution is artificial.

A preliminary integration run exposed a redundant pinpoint escalation and an incorrect retrieval denominator, which were corrected before the recorded final run. The same fixed split was reused; this is a transparent engineering check, not a blind preregistered evaluation. Fixed context cues and the 0.3 candidate-display threshold were not optimized against test F1. A genuinely independent external benchmark remains necessary.

An unknown alias for a real case produces unresolved even if another alias would succeed. A PDF with missing or ambiguous numbered text can prevent a valid quote from matching. Similar but substantively different propositions remain outside the tool's competence. Short or repeated quotes may provide insufficient context for qualification warnings. Accurate text can also reproduce counsel's argument, a lower court's reasoning, or an overruled authority; none is settled by contiguous matching. Consequently, overall triage F1 must not be marketed as the accuracy of legal advice.

Tests cover boundaries, negation and number changes, duplicate aliases, ambiguous identity, unknown metadata, repeated quotes at valid pinpoints, malformed inputs, and typed evidence. CLI and HTTP tests use isolated temporary data. The browser exercises the real engine rather than scripted mock responses. Per-case predictions and metrics are stored as machine-readable files; PDF measurement macros are generated from them to avoid manual transcription.

The repository includes source code, tests, an attributed retained corpus, manifest, experiment commands, demonstration instructions, and editable report source. Raw PDFs and model weights are not committed. Experiments can use the retained corpus without network acquisition; reacquisition depends on upstream availability. Team contributions require truthful attestation rather than invented equal assignments. The rubric requires every team member to explain the entire project and its limitations during the viva, so the demonstration guide includes both successful and unresolved checks.

\section{Conclusion}
This project implements a bounded Indian citation-integrity workflow that resolves metadata-backed authorities, checks quotations and mechanical metadata claims, returns original paragraph evidence, and distinguishes discrepancies from unresolved or review-required cases. The controlled held-out experiment gives \FullFone{} F1 for the complete pipeline; simpler baselines expose the limits of checking only existence or text. These results support the implemented mechanical behavior, not a claim of legal understanding or equivalent performance on real filings. Corpus incompleteness, PDF ambiguity, English-only processing, automatic labels, and cue-based context warnings remain limitations. Next steps are an independently annotated natural-error benchmark, broader source coverage with verified aliases, opinion-aware extraction, and expert evaluation of proposition support. Until then, the system is an assistive evidence checker rather than a replacement for professional legal review.

\begin{thebibliography}{20}
\bibitem{ref1} Matthew Dahl, Varun Magesh, Mirac Suzgun, and Daniel E. Ho, ``Large Legal Fictions: Profiling Legal Hallucinations in Large Language Models,'' \emph{Journal of Legal Analysis}, 2024, doi: 10.1093/jla/laae003.
\bibitem{ref2} Varun Magesh, Faiz Surani, Matthew Dahl, Mirac Suzgun, Christopher D. Manning, and Daniel E. Ho, ``Hallucination-Free? Assessing the Reliability of Leading AI Legal Research Tools,'' \emph{Journal of Empirical Legal Studies}, 2025, doi: 10.1111/jels.12413.
\bibitem{ref3} Patrick Lewis et al., ``Retrieval-Augmented Generation for Knowledge-Intensive NLP Tasks,'' \emph{NeurIPS}, 2020.
\bibitem{ref4} Tianyu Gao et al., ``Enabling Large Language Models to Generate Text with Citations,'' \emph{EMNLP}, 2023, doi: 10.18653/v1/2023.emnlp-main.398.
\bibitem{ref5} James Thorne, Andreas Vlachos, Christos Christodoulopoulos, and Arpit Mittal, ``FEVER: a Large-scale Dataset for Fact Extraction and VERification,'' \emph{NAACL-HLT}, 2018, doi: 10.18653/v1/N18-1074.
\bibitem{ref6} David Wadden et al., ``Fact or Fiction: Verifying Scientific Claims,'' \emph{EMNLP}, 2020, doi: 10.18653/v1/2020.emnlp-main.609.
\bibitem{ref7} Wojciech Kryscinski et al., ``Evaluating the Factual Consistency of Abstractive Text Summarization,'' \emph{EMNLP}, 2020, doi: 10.18653/v1/2020.emnlp-main.750.
\bibitem{ref8} Potsawee Manakul, Adian Liusie, and Mark Gales, ``SelfCheckGPT: Zero-Resource Black-Box Hallucination Detection for Generative Large Language Models,'' \emph{EMNLP}, 2023, doi: 10.18653/v1/2023.emnlp-main.557.
\bibitem{ref9} Neel Guha et al., ``LegalBench: A Collaboratively Built Benchmark for Measuring Legal Reasoning in Large Language Models,'' \emph{NeurIPS Datasets and Benchmarks}, 2023.
\bibitem{ref10} Ilias Chalkidis, Abhik Jana, Dirk Hartung, Michael Bommarito, Ion Androutsopoulos, Daniel Katz, and Nikolaos Aletras, ``LexGLUE: A Benchmark Dataset for Legal Language Understanding in English,'' \emph{ACL}, 2022, doi: 10.18653/v1/2022.acl-long.297.
\bibitem{ref11} Vijit Malik, Rishabh Sanjay, Shubham Kumar Nigam, Kripabandhu Ghosh, Shouvik Kumar Guha, Arnab Bhattacharya, and Ashutosh Modi, ``ILDC for CJPE: Indian Legal Documents Corpus for Court Judgment Prediction and Explanation,'' \emph{ACL-IJCNLP}, 2021, doi: 10.18653/v1/2021.acl-long.313.
\bibitem{ref12} Jacob Devlin et al., ``BERT: Pre-training of Deep Bidirectional Transformers for Language Understanding,'' \emph{NAACL-HLT}, 2019, doi: 10.18653/v1/N19-1423.
\bibitem{ref13} Nils Reimers and Iryna Gurevych, ``Sentence-BERT: Sentence Embeddings using Siamese BERT-Networks,'' \emph{EMNLP-IJCNLP}, 2019, doi: 10.18653/v1/D19-1410.
\bibitem{ref14} Ilias Chalkidis et al., ``LEGAL-BERT: The Muppets straight out of Law School,'' \emph{Findings of EMNLP}, 2020, doi: 10.18653/v1/2020.findings-emnlp.261.
\bibitem{ref15} Stephen Robertson and Hugo Zaragoza, ``The Probabilistic Relevance Framework: BM25 and Beyond,'' \emph{Foundations and Trends in Information Retrieval}, 2009, doi: 10.1561/1500000019.
\bibitem{ref16} Vladimir Karpukhin et al., ``Dense Passage Retrieval for Open-Domain Question Answering,'' \emph{EMNLP}, 2020, doi: 10.18653/v1/2020.emnlp-main.550.
\bibitem{ref17} Nandan Thakur et al., ``BEIR: A Heterogeneous Benchmark for Zero-shot Evaluation of Information Retrieval Models,'' \emph{NeurIPS Datasets and Benchmarks}, 2021, doi: 10.48550/arXiv.2104.08663.
\bibitem{ref18} Ashish Vaswani et al., ``Attention Is All You Need,'' \emph{NeurIPS}, 2017.
\bibitem{ref19} Or Honovich et al., ``TRUE: Re-evaluating Factual Consistency Evaluation,'' \emph{NAACL}, 2022, doi: 10.18653/v1/2022.naacl-main.287.
\bibitem{dataset} AWS Open Data Registry; Vanga, ``Indian Supreme Court Judgments,'' \emph{Registry of Open Data on AWS}, 2025. [Online]. Available: \url{https://registry.opendata.aws/indian-supreme-court-judgments/}.
\end{thebibliography}

\end{document}
